\chapter{Muscle Pain}

\section*{Definition}
Pain discomfort or tenderness in skeletal muscles often described aching soreness cramping or stiffness.

\section*{When to See a Doctor}
See your doctor if:
\begin{itemize}
  \item Muscle pain is severe, unexplained, or getting worse
\end{itemize}

\section*{How It Develops}
Can be acute from exertion injury infection or chronic from inflammatory autoimmune medication metabolic conditions; myalgia differs from myositis (muscle inflammation).

\section*{Causes}
\begin{itemize}
  \item Viral infections (influenza COVID-19)
  \item Overuse/exercise
  \item Statin medications
  \item Polymyalgia rheumatica morning stiffness shoulder hips elderly
  \item Electrolyte abnormalities (hypokalemia hypocalcemia)
  \item Fibromyalgia
\end{itemize}

\section*{Related Symptoms}
\begin{itemize}
  \item Muscle weakness
  \item Fatigue
  \item Stiffness
  \item Swelling at site
  \item Feeling of muscle nodules
\end{itemize}


\section*{Medical Evaluation}
\begin{itemize}
  \item History and examination assess strength, tenderness, distribution, exertion, medication exposure, infection, and systemic symptoms.
  \item Tests may include creatine kinase, kidney function, electrolytes, thyroid studies, inflammatory markers, or imaging according to the suspected cause.
\end{itemize}

\section*{Treatment Options}
\begin{itemize}
  \item Treatment addresses the cause, such as stopping an offending medicine, correcting an electrolyte disorder, treating infection, or managing inflammatory or autoimmune disease.
  \item Rhabdomyolysis, progressive weakness, or severe inflammatory myopathy requires urgent or specialist-directed care.
\end{itemize}

\section*{Self-Care and Home Management}
\begin{itemize}
  \item Reduce aggravating activity temporarily but resume gentle movement as tolerated.
  \item Fluids, sleep, heat or ice, and an appropriate over-the-counter analgesic may help uncomplicated soreness.
  \item Seek urgent care for marked weakness, dark or cola-coloured urine, severe swelling, high fever, or pain after heat exposure or major exertion.
\end{itemize}

\section*{References}

\begin{thebibliography}{99}
\bibitem{muscle_pain:harrison-myalg} Longo DL, et al. \emph{Harrison's Principles of Internal Medicine}. 21st ed. McGraw-Hill; 2022. Chapter 435: Myalgia.
\bibitem{muscle_pain:uptodate-myalg} Feld JJ, et al. Approach to the patient with myalgia. In: UpToDate, Post TW (Ed), UpToDate, Waltham, MA; 2025.
\bibitem{muscle_pain:fibro-guideline} Fitzcharles MA, et al. 2023 Canadian guideline for the diagnosis and management of fibromyalgia. \emph{Pain}. 2024;165(2):267-283.
\bibitem{muscle_pain:goldstein} Goldstein MR, et al. Statin-associated muscle symptoms: a clinical review. \emph{J Am Coll Cardiol}. 2023;81(15):1489-1502.
\bibitem{muscle_pain:diamanti} Diamanti AP, et al. Polymyalgia rheumatica: an update on diagnosis and treatment. \emph{Lancet Rheumatol}. 2024;6(2):e96-e107.
\bibitem{muscle_pain:bialas} Bialas AJ, et al. Viral myalgia: mechanisms and management. \emph{Clin Infect Dis}. 2023;76(10):e1820-e1830.
\end{thebibliography}
